\section{Déclaration du type}
\begin{algorithme}

\begin{enregistrement}{ArbreHuffman}
    \champEnregistrement{statistiques}{Statistiques}
    \champEnregistrement{filsGauche}{\typePointeur{ArbreHuffman}}
    \champEnregistrement{filsDroit}{\typePointeur{ArbreHuffman}}
\end{enregistrement}
\end{algorithme}

 
\subsubsection{arbreDeHuffman}

\begin{algorithme}

    \fonction{arbreDeHuffman}{stat : Statistiques}{\typePointeur{ArbreHuffman}}{}
        {arbre : \typePointeur{ArbreHuffman}}
        {
        \sialors{non(estVide(arbre))}
        {
        \affecter{\champPointeur{arbre}{statistiques}}{stat}
        \affecter{\champPointeur{arbre}{filsGauche}}{NIL}
        \affecter{\champPointeur{arbre}{filsDroit}}{NIL}
        }
        \retourner{arbre}        
        }
\end{algorithme}

\subsubsection{fixerStatistique}

\begin{algorithme}

    \procedure{fixerStatistique}{\paramEntreeSortie{arbre : \typePointeur{ArbreHuffman}}\paramEntree{stat : Statistiques}}{}
    {}
    {
    \affecter{\champPointeur{arbre}{statistiques}}{stat}
    }


\end{algorithme}

\subsubsection{obtenirStat}

\begin{algorithme}

    \fonction{obtenirStat}{arbre : ArbreHuffman}{Statistiques}{}
    {}
    {
    \retourner{\champ{arbre}{statistiques}}
    }


\end{algorithme}

\subsubsection{obtenirPonderation}

\begin{algorithme}
    \fonction{obtenirPonderation}{arbre : ArbreHuffman}{\naturel}{}{}
    {
        \retourner{\champ{arbre}\champ{statistiques}{ponderation}}

    }
\end{algorithme}

\subsubsection{obtenirCaracter}

\begin{algorithme}
    \fonction{obtenirCaracter}{arbre : ArbreHuffman}{\caractere}{}{}
    {
        \retourner{\champ{arbre}\champ{statistiques}{caractere}}

    }
\end{algorithme}

\subsubsection{fusionner}

\begin{algorithme}

    \fonction{fusionner}{arbre1 , arbre2 : \typePointeur{ArbreHuffman}}{\typePointeur{ArbreHuffman}}{non (estVide(arbre1) et estVide(arbre2))}
    {res : \typePointeur{ArbreHuffman}, stat : Statistiques}
    {   
        \affecter{\champ{stat}{caractere}}{' '}
        \affecter{\champ{stat}{ponderation}}{obtenirPonderation(\typePointeur{arbre1})+obtenirPonderation(\typePointeur{arbre2})}
        \affecter{res}{arbreDeHuffman(stat)}
        \sialorssinon{obtenirPonderation(\typePointeur{arbre1})$<$obtenirPonderation(\typePointeur{arbre2})}
        {fixerFilsGauche(res,arbre1)\\
        fixerFilsDroit(res,arbre2)}
        {fixerFilsGauche(res,arbre2)\\
        fixerFilsDroit(res,arbre1)}
        \retourner{res}
        

    }
\end{algorithme}


\subsubsection{fixerFilsDroit}

\begin{algorithme}
    \procedure{fixerFilsDroit}{\paramEntreeSortie{arbre : \typePointeur{ArbreHuffman}} \paramEntree{droit : \typePointeur{ArbreHuffman}}}{}{}
        {
            \sialors{non(estVide(arbre))}
            {
            \affecter{\champPointeur{arbre}{filsDroit}}{droit}
            }
        }

    

\end{algorithme}

\subsubsection{fixerFilsGauche}

\begin{algorithme}
    \procedure{fixerFilsGauche}{\paramEntreeSortie{arbre : \typePointeur{ArbreHuffman}} \paramEntree{gauche : \typePointeur{ArbreHuffman}}}{}{}
        {
            \sialors{non(estVide(arbre))}{
            \affecter{\champPointeur{arbre}{filsGauche}}{gauche}}
        }
\end{algorithme}

\subsubsection{estVide}

\begin{algorithme}
    \fonction{estVide}{arbre : \typePointeur{ArbreHuffman}}{\booleen}{}{}
    {
        \retourner{((\champPointeur{arbre}\champ{statistiques}{caractere} = NULL) et (\champ{arbre}\champ{statistique}{ponderation} = NULL))}
    }
\end{algorithme}

\subsubsection{estFeuille}

\begin{algorithme}
    \fonction{estFeuille}{arbre : ArbreHuffman}{\booleen}{}{}
    {

    \retourner{(estVide(obtenirFilsGauche(@arbre)) et estVide(obtenirFilsDroit(&arbre)))}
    }
\end{algorithme}


\subsubsection{obtenirFilsGauche}

\begin{algorithme}
    \fonction{obtenirFilsGauche}{arbre : \typePointeur{ArbreHuffman}}{\typePointeur{ArbreHuffman}}{non estFeuille(arbre)}{}
    {
    \retourner{\champPointeur{arbre}{filsGauche}}
    }
\end{algorithme}

\subsubsection{obtenirFilsDroit}

\begin{algorithme}
    \fonction{obtenirFilsDroit}{arbre : \typePointeur{ArbreHuffman}}{\typePointeur{ArbreHuffman}}{non estFeuille(arbre)}{}
    {
    \retourner{\champPointeur{arbre}{filsDroit}}
    }
\end{algorithme}

\subsubsection{estPresent}

\begin{algorithme}
    \fonction{estPresent}{arbre : \typePointeur{ArbreHuffman}, caractere : \caractere}{\booleen}{}{}
    {   
        \sialors(non(estVide(arbre)))
        {
        \sialorssinon{\champPointeur{arbre}\champ{statistiques}{caractere} = caractere}
        {\retourner{Vrai}}
        {\retourner{estPresent(obtenirFilsDroit(arbre)) ou estPresent(obtenirFilsGauche(arbre))}}
        }
        }
\end{algorithme}
